% ===============================
%           PAQUETES
% ===============================

\documentclass[8pt]{article}

\usepackage[spanish,es-noshorthands]{babel}
\usepackage[utf8]{inputenc}
\usepackage{xcolor}
\usepackage{amsmath, amssymb}
\usepackage{graphicx}
\usepackage{fancyhdr}
\usepackage{tcolorbox}
\usepackage{titlesec}
\usepackage{hyperref}
\usepackage{tikz}
\usepackage{multicol}
\tcbuselibrary{breakable}
\usepackage{pgfplots}
\pgfplotsset{compat=1.18}
\usepackage{colortbl}
\usetikzlibrary{arrows.meta, shapes, positioning, calc, patterns}

\usepackage{geometry}
\geometry{margin=2cm}

% ===============================
%           COLORES
% ===============================

\definecolor{azuloscuro}{RGB}{0, 51, 102}
\definecolor{azulmedio}{RGB}{42, 82, 152}
\definecolor{verdecorrecto}{RGB}{40, 167, 69}
\definecolor{fondogris}{RGB}{248, 249, 250}
\definecolor{fondoverde}{RGB}{235, 247, 238}
\definecolor{amarillonota}{RGB}{255, 250, 205}
\definecolor{fondorojo}{RGB}{248, 215, 218}

% ===============================
%     ESTILO DE ENCABEZADO
% ===============================

\newcommand{\materia}{Área 1: Físico-Matemáticas}
\newcommand{\tema}{}

\pagestyle{fancy}
\fancyhf{}
\lhead{\tema}
\rhead{\materia}
\cfoot{\thepage}

% ===============================
%     ESTILO DE SECCIONES
% ===============================

\newtcolorbox{seccion}{
    colback=fondogris,
    colframe=azulmedio,
    boxrule=0pt,
    leftrule=4pt,
    arc=3pt,
    breakable
}

% ===============================
%     CAJA PARA DEFINICIONES
% ===============================

\newtcolorbox{definicion}{
colback=fondogris,
colframe=azuloscuro,
arc=4pt,
boxrule=1pt
}

% ===============================
%     CAJA PARA ECUACIONES
% ===============================

\newtcolorbox{ecuacion}{
    colback=fondoverde,
    colframe=verdecorrecto,
    boxrule=0pt,
    leftrule=3pt,
    arc=3pt,
    breakable
}

% ===========================
%     NOTAS IMPORTANTES
% ===========================

\newtcolorbox{nota}{
    colback=amarillonota,
    colframe=azuloscuro,
    boxrule=1pt,
    arc=3pt,
    breakable
}

% ======================
%     INSTRUCCIONES
% ======================

\newtcolorbox{instruccionbox}{
    colback=amarillonota,
    colframe=azuloscuro,
    boxrule=1pt,
    arc=3pt,
    breakable
}

% ==================
%     PREGUNTAS
% ==================

\newtcolorbox{preguntabox}{
    colback=fondogris,
    colframe=azulmedio,
    boxrule=0pt,
    leftrule=3pt,
    arc=3pt,
    breakable
}

% ===================
%     RESPUESTAS
% ===================

\newtcolorbox{respuestabox}{
    colback=fondoverde,
    colframe=verdecorrecto,
    boxrule=0pt,
    leftrule=4pt,
    arc=3pt,
    breakable
}


